%% =============================================================================
%% 數學九章 (Shuxue Jiuzhang) — Fascicles 5–6
%% Author: Qin Jiushao (秦九韶), Song Dynasty
%% Source: Siku Quanshu (四庫全書) Edition
%% Typeset with XeLaTeX
%% =============================================================================
\documentclass[11pt,a4paper]{article}

%% -------- Core Packages --------
\usepackage{fontspec}
\usepackage{polyglossia}
\usepackage{amsmath,amssymb}
\usepackage{geometry}
\usepackage{graphicx}
\usepackage{xcolor}
\usepackage{titlesec}
\usepackage{fancyhdr}
\usepackage{enumitem}
\usepackage{array,booktabs,longtable,multirow}
\usepackage{hyperref}
\usepackage{tocloft}
\usepackage{indentfirst}
\usepackage{xeCJK}

%% -------- Page Geometry --------
\geometry{
  paperwidth=210mm,paperheight=297mm,
  left=25mm,right=25mm,top=30mm,bottom=30mm,
  headheight=14pt,footskip=30pt
}

%% -------- Language Setup --------
\setmainlanguage{english}
\setotherlanguage{chinese}

%% -------- Font Configuration --------
\setmainfont{LinLibertine_R.otf}[
  BoldFont=LinLibertine_RB.otf,
  ItalicFont=LinLibertine_RI.otf,
  BoldItalicFont=LinLibertine_RBI.otf,
  Renderer=Harfbuzz
]
\setsansfont{LinBiolinum_R.otf}[
  BoldFont=LinBiolinum_RB.otf,
  ItalicFont=LinBiolinum_RI.otf,
  Renderer=Harfbuzz
]

%% Chinese Fonts
\setCJKmainfont{Noto Sans CJK SC}
\setCJKsansfont{Noto Sans CJK SC}
\setCJKmonofont{Noto Sans CJK SC}

%% -------- Colors --------
\definecolor{titlecolor}{RGB}{45,55,72}
\definecolor{sectioncolor}{RGB}{55,65,81}
\definecolor{notecolor}{RGB}{120,53,15}
\definecolor{rulecolor}{RGB}{180,160,140}
\definecolor{problemcolor}{RGB}{80,40,20}
\definecolor{answercolor}{RGB}{20,80,40}
\definecolor{methodcolor}{RGB}{20,40,80}

%% -------- Section Formatting --------
\titleformat{\section}
  {\normalfont\Large\bfseries\color{sectioncolor}}
  {\thesection}{1em}{}
\titleformat{\subsection}
  {\normalfont\large\bfseries\color{sectioncolor}}
  {\thesubsection}{1em}{}
\titleformat{\subsubsection}
  {\normalfont\normalsize\bfseries\color{problemcolor}}
  {\thesubsubsection}{1em}{}

%% -------- Header/Footer --------
\pagestyle{fancy}
\fancyhf{}
\fancyhead[L]{\small\textsc{數學九章 — Fascicles 5–6}}
\fancyhead[R]{\small\thepage}
\renewcommand{\headrulewidth}{0.4pt}
\renewcommand{\footrulewidth}{0pt}

%% -------- Custom Commands --------
\newcommand{\longline}{\noindent\rule{\textwidth}{0.4pt}}
\newcommand{\shortline}{\noindent\rule{0.5\textwidth}{0.4pt}\par}
\newcommand{\cnote}[1]{\textcolor{notecolor}{\textit{#1}}}

%% Problem structure commands
\newcommand{\Question}[1]{\par\noindent\textbf{\color{problemcolor}問：}#1\par}
\newcommand{\Answer}[1]{\par\noindent\textbf{\color{answercolor}答曰：}#1\par}
\newcommand{\Method}[1]{\par\noindent\textbf{\color{methodcolor}術曰：}#1\par}
\newcommand{\Working}[1]{\par\noindent\textbf{草曰：}#1\par}
\newcommand{\EditorialNote}[1]{\par\noindent\textit{［按］#1}\par}

%% -------- Hyperref Setup --------
\hypersetup{
  colorlinks=true,
  linkcolor=sectioncolor,
  citecolor=sectioncolor,
  urlcolor=sectioncolor,
  pdfencoding=auto,
  pdftitle={數學九章 — Fascicles 5–6},
  pdfauthor={Qin Jiushao (秦九韶)},
  pdfsubject={Chinese Mathematics}
}

%% -------- TOC Depth --------
\setcounter{tocdepth}{3}
\setcounter{secnumdepth}{3}

%% Reduce paragraph spacing for dense Chinese text
\setlength{\parskip}{0.5ex plus 0.2ex minus 0.1ex}
\setlength{\parindent}{2em}

%% =============================================================================
\begin{document}

%% -------- Title Page --------
\begin{titlepage}
\centering
\vspace*{3cm}

{\fontsize{36}{44}\selectfont\color{titlecolor}\CJKfamily{}\textbf{數學九章}}

\vspace{1.5cm}

{\fontsize{18}{22}\selectfont\color{sectioncolor} Fascicles Five and Six}

\vspace{2cm}

{\Large\color{sectioncolor} 卷五：\textbf{賦役} \quad|\quad 卷六：\textbf{錢穀}}

\vspace{3cm}

{\large 宋\quad 秦九韶\quad 撰}

\vspace{0.5cm}

{\normalsize\textit{Qin Jiushao (c.\ 1202–1261)}}

\vspace{2cm}

{\color{rulecolor}\rule{0.6\textwidth}{0.4pt}}

\vspace{1.5cm}

{\small\color{sectioncolor}
\textit{Siku Quanshu} (四庫全書) Edition\\[0.5em]
Digital Typesetting\quad|\quad Modern Edition}

\vspace{3cm}

{\small\color{notecolor}
\textit{This document preserves the traditional Chinese mathematical\\
text structure of 問 (Problem), 答 (Answer), 術 (Method), and 草 (Working).}}

\vfill
{\small\color{sectioncolor} Compiled with XeLaTeX}
\end{titlepage}

%% -------- Table of Contents --------
\tableofcontents
\newpage


%% =============================================================================
%% FASCICLE 5A — 賦役 (Part 1)
%% =============================================================================
\newpage
\section*{\color{titlecolor}卷五上\quad 賦役}
\addcontentsline{toc}{section}{卷五上\quad 賦役}

\noindent\textbf{按：}賦役一章，即古九章中差分粟米諸法。但從來算書設題之數未有如此卷之多者。如首題答數至一百七十五條，毎條歩算之數至十餘位，而得數皆無不合，亦可謂能廣其用矣。且諸問皆足以明貴賤廩税及交質變易之同異，使公私各得其平，誠治民之要務也。但題語多晦，術草中間有與所問不相應者，亦於各條下指出，俾觀者無惑焉。

\longline

\subsection*{\color{sectioncolor}復邑修賦}
\addcontentsline{toc}{subsection}{復邑修賦}

\Question{有海坍縣地，今有復漲，歳乆鄉井再成，申諸剏邑。稱土排到六鄉，以附郭為甲，最逺為已，各有田九等，開具下項：

\begin{longtable}{p{0.95\textwidth}}
\toprule
\textbf{甲鄉}共計田十四萬一百九十三畝三角一十二步\\
\midrule
\endhead
上等：上田五千六百七十八畝一角四十八歩；中田四千八百九十二畝三十歩；下田六千六百二十一畝五十四歩\\
中等：上田八千二百二十五畝二十四歩；中田一萬三十五畝六歩；下田一萬六千五百三十畝\\
下等：上田二萬一千九十畝二十四歩；中田三萬二千六十畝三歩；下田三萬五千六十一畝三角三歩\\
\bottomrule
\end{longtable}

\begin{longtable}{p{0.95\textwidth}}
\toprule
\textbf{乙鄉}田共計八萬四千一十畝二角二歩\\
\midrule
\endhead
上等：上田六千七百八十九畝一角三十六歩；中田五千九百八十七畝二角；下田八千一十畝三角三歩\\
中等：上田七千五百四十一畝；中田九千一百二十一畝二角一十二歩；下田一萬九千六十歩\\
下等：上田一萬八千三十七畝一角六歩；中田九千四百五十六畝三角五十九歩；下田無\\
\bottomrule
\end{longtable}

\begin{longtable}{p{0.95\textwidth}}
\toprule
\textbf{丙鄉}田共計一十二萬九百三十五畝五十八歩五分\\
\midrule
\endhead
上等：上田四千八百六十八畝二角三歩；中田五千九百七十九畝三角六歩；下田六千八百八十八畝二角六歩\\
中等：上田七千九百八十四畝一角；中田一萬四千五十六畝一十二歩；下田二萬三千三百三十三畝一十二歩\\
下等：上田二萬七千七百五十五畝一十六歩五分；中田無；下田三萬七十畝三歩\\
\bottomrule
\end{longtable}

\begin{longtable}{p{0.95\textwidth}}
\toprule
\textbf{丁鄉}田共計八萬九千六十六畝二歩三分\\
\midrule
\endhead
上等：上田一萬一千一百二畝一歩；中田九千八百七十六畝一角；下田八千七百六十五畝一角三十歩\\
中等：上田七千五百三十九畝三十四歩三分；中田無；下田一萬二千九百八十七畝四十二歩\\
下等：上田無；中田五千四百三十二畝一角六歩；下田三萬三千三百六十三畝三角九歩\\
\bottomrule
\end{longtable}

\begin{longtable}{p{0.95\textwidth}}
\toprule
\textbf{戊鄉}田共計二十萬四千四百七十四畝一角二十四歩四分\\
\midrule
\endhead
上等：上田二萬四千六百三十二畝三十九歩；中田一萬三千五百二十一畝二十七歩；下田九千九百八十八畝三角三歩\\
中等：上田八千八百七十七畝五十六歩四分；中田一萬一千三百三十三畝三角；下田二萬七千六十七畝\\
下等：上田一萬九千八百七十六畝三角六歩；中等田七萬九千一百三十五畝三角四十三歩；下田一萬四十一畝二角三十歩\\
\bottomrule
\end{longtable}

\begin{longtable}{p{0.95\textwidth}}
\toprule
\textbf{已鄉}田共計一十五萬八千四百六十畝三角一十八歩二分\\
\midrule
\endhead
上等：上田無；中田七千七百八十八畝三角五十一歩；下田無\\
中等：上田九千九百九十九畝二角三歩；中田一萬八百三十六畝五十六歩；下田無\\
下等：上田三萬二千八十九畝一角四十五歩六分；中田四萬三千六百七十八畝二角五十七歩；下田五萬四千六十七畝三角四十五歩六分\\
\bottomrule
\end{longtable}

照得昨來本縣原科苗米一十萬三千五百六十七石八斗四升四合三勺，和買一萬三千四百九十八匹一丈七尺三寸七分六釐，夏税九千八百七十六匹三丈二尺六寸五分八釐。其六鄉田係三色，甲為上，乙丙為次，丁戊己又為次。先令官物為三差，使上比中、中比下皆十分外差一，次令各鄉九等皆於十分内差一，抛科用合租額。其乙鄉田最肥，次丁，次甲，次丙，次己，次戊。欲知三色九等毎畝等則及共科數各鄉幾何？
}

\EditorialNote{題内言田有三色，甲為上云云，又言乙鄉田最肥云云，語若相左。葢前言三色者六鄉共科之等，後言田肥者每畝所出之異也。若易田係三色為科有三差，則語意明矣。}

\Answer{
\textbf{甲鄉：}上等上田苗三斗二升三合二勺，和一尺六寸九分，税一尺二寸三分；
中田苗二斗九升九勺，和一尺五寸二分，税一尺一寸一分；
下田苗二斗五升八合六勺，和一尺三寸五分，税九寸九分。

中等上田苗二斗二升六合二勺，和一尺一寸八分，税八寸六分；
中田苗一斗九升三合九勺，和一尺一分，税七寸四分；
下田苗一斗六升一合六勺，和八寸四分，税六寸二分。

下等上田苗一斗二升九合三勺，和六寸七分，税四寸九分；
中田苗九升六合九勺，和五寸一分，税三寸七分；
下田苗六升四合六勺，和三寸四分，税二寸五分。

\textbf{乙鄉：}上等上田苗三斗六升三合三勺，和一尺八寸九分，税一尺三寸九分；
中田苗三斗二升六合九勺，和一尺七寸，税一尺二寸五分；
下田苗二斗九升六勺，和一尺五寸二分，税一尺一寸一分。

中等上田苗二斗五升四合三勺，和一尺三寸三分，税九寸七分；
中田苗二斗一升八合，和一尺一寸四分，税八寸三分；
下田苗一斗八升一合六勺，和九寸五分，税六寸九分。

下等上田苗一斗四升五合三勺，和七寸六分，税五寸五分；
中田苗一斗九合，和五寸七分，税四寸二分；
下田苗七升二合七勺，和三寸八分，税二寸八分。

\textbf{丙鄉：}上等上田苗三斗三合五勺，和一尺五寸八分，税一尺一寸六分；
中田苗二斗七升三合一勺，和一尺四寸二分，税一尺四分；
下田苗二斗四升二合八勺，和一尺二寸七分，税九寸三分。

中等上田苗二斗一升二合四勺，和一尺一寸一分，税八寸一分；
中田苗一斗八升二合一勺，和九寸五分，税六寸九分；
下田苗一斗五升一合七勺，和七寸九分，税五寸八分。

下等上田苗一斗二升一合四勺，和六寸三分，税四寸六分；
中田苗九升一合，和四寸七分，税三寸五分；
下田苗六升七勺，和三寸二分，税二寸三分。

\textbf{丁鄉：}上等上田苗三斗四升三合二勺，和一尺七寸九分，税一尺三寸一分；
中田苗三斗八合九勺，和一尺六寸一分，税一尺一寸八分；
下田苗二斗七升四合六勺，和一尺四寸三分，税一尺五分。

中等上田苗二斗四升二勺，和一尺二寸五分，税九寸二分；
中田苗二斗五合九勺，和一尺七分，税七寸九分；
下田苗一斗七升一合六勺，和八寸九分，税六寸五分。

下等上田苗一斗三升七合三勺，和七寸二分，税五寸二分；
中田苗一斗三合，和五寸四分，税三寸九分；
下田苗六升八合六勺，和三寸六分，税二寸六分。

\textbf{戊鄉：}上等上田苗一斗五升三合八勺，和八寸，税五寸九分；
中田苗一斗三升八合四勺，和七寸二分，税五寸三分；
下田苗一斗二升三合一勺，和六寸四分，税四寸七分。

中等上田苗一斗七合七勺，和五寸六分，税四寸一分；
中田苗九升二合三勺，和四寸八分，税三寸五分；
下田苗七升六合九勺，和四寸，税二寸九分。

下等上田苗六升一合五勺，和三寸二分，税二寸三分；
中田苗四升六合一勺，和二寸四分，税一寸八分；
下田苗三升八勺，和一寸六分，税一寸二分。

\textbf{已鄉：}上等上田苗二斗八升二合二勺，和一尺四寸七分，税一尺八分；
中田苗二斗五升三合九勺，和一尺三寸二分，税九寸七分；
下田苗二斗二升五合七勺，和一尺一寸八分，税八寸六分。

中等上田苗一斗九升七合五勺，和一尺三分，税七寸五分；
中田苗一斗六升九合三勺，和八寸八分，税六寸五分；
下田苗一斗四升一合一勺，和七寸四分，税五寸四分。

下等上田苗一斗一升二合九勺，和五寸九分，税四寸三分；
中田苗八升四合六勺，和四寸四分，税三寸二分；
下田苗五升六合四勺，和二寸九分，税二寸二分。

\textbf{各鄉共科數：}
甲鄉苗米一萬九千五百五十石二斗四升八合三勺；
和買一萬一百九十二丈二尺六寸五分六釐；
夏税七千四百五十七丈六尺八寸九分八釐。

乙鄉苗米一萬七千七百七十二石九斗五升三合；
和買九千二百六十五丈六尺九寸六分；
夏税六千七百七十九丈七尺一寸八分。

丙鄉苗米一萬七千七百七十二石九斗五升三合；
和買九千二百六十五丈六尺九寸六分；
夏税六千七百七十九丈七尺一寸八分。

丁鄉苗米一萬六千一百五十七石二斗三升；
和買八千四百二十三丈三尺六寸；
夏税六千一百六十三丈三尺八寸。

戊鄉苗米一萬六千一百五十七石二斗三升；
和買八千四百二十三丈三尺六寸；
夏税六千一百六十三丈三尺八寸。

已鄉苗米一萬六千一百五十七石二斗三升；
和買八千四百二十三丈三尺六寸；
夏税六千一百六十三丈三尺八寸。
}

\Method{以衰分求之。先列本縣色位，自下錐行列之，又以鄉數對列而乗之，副併為法，以除官物數，得一分之率。以率數乗未併者，各得諸鄉之數。次列各鄉等位，自上等置十分，每以内分錐行九折之，至九等止，又各以畝歩乗之，副併為鄉法，以除諸各鄉所得官物數，所得為一分之率，以乗未併者，各得毎畝税色。}

\Working{列本縣色位三，自下色例十分中十一分，上比中身外加一，上得一百二十一，中得一百一十，下得一百，為上中下三率，列之右行。乃列甲一對上率，乙丙共二對中率，丁戊己共三對下率，列左行。乃以左行列數各相對乗右行率數，其上得一百二十一，其十得二百二十，其下得三百，乃副置而併之，得六百四十一為法。置元科苗米一萬三千五百六十七石八斗四升四合三勺為寔，以法除之，得一百六十一石五斗七升二合三勺為一分之率。以未併下率一百乗率，得一萬六千一百五十七石二斗三升為丁戊己三鄉各科數。以於身下加一，得一萬七千七百七十二石九斗五升三合為乙丙二鄉各科數。又於身下加一，得一萬九千五百五十石二斗四升八合三勺為甲合科數。求和買亦用六百四十一為法，置元科和買一萬三千四百九十八匹，以匹法四丈通之，得五萬三千九百九十二丈，内零一丈七尺三寸七分六釐，得五萬三千九百九十三丈七尺三寸七分六釐為寔，以法除之，得八十四丈二尺三寸三分六釐為一分之率。亦以未併下率一百乗之，得八千四百二十三丈三尺六寸為丁戊己三鄉各科數。次於身下加一，得九千二百六十五丈六尺九寸六分為乙丙二鄉各科數。又於身下加一，得一萬一百九十二丈二尺六寸五分六釐為甲鄉合科數。求夏税亦置六百四十一為法，置元科九千八百七十六匹，以匹法四丈通之，得三萬九千五百四丈，内零三丈二尺六寸五分八釐，得三萬九千五百七丈二尺六寸五分八釐為寔，以法除之，得六十一丈六尺三寸三分八釐為一分率。亦以未併下率一百乗之，得六千一百六十三丈三尺八寸為丁戊己三鄉各科數。次於身下加一，得六千七百七十九丈七尺一寸八分为乙丙二鄉各科數。又於身下加一，得七千四百五十七丈六尺八寸九分八釐為甲鄉數。}

\newpage
\subsection*{\color{sectioncolor}圍田租畝}
\addcontentsline{toc}{subsection}{圍田租畝}

\Question{有興復圍田已成，共計三千二十一頃五十一畝一十五歩，分三等。其上等毎畆起租六斗，中等四斗五升，下等四斗。中田多上田弱半，不及下田太半。欲知三色田畝及各租幾何？}

\EditorialNote{題言中田多上田弱半，不及下田太半。蓋以弱半為三分之一，太半為三分之二也。多上田弱半者，即比上田多三分之一也。不及下田太半者，即比下田少三分之二也。是上田加三分之一為中田，即下田三分之一也。轉言之，則中田為下田三分之一，上田為中田四分之三也。}

\Answer{上田四百七十七頃八畝一十五歩，米二萬八千六百二十四石八斗三升七合五勺；
中田六百三十六頃一十畝三角，米二萬八千六百二十四石八斗三升七合五勺；
下田一千九百八頃三十二畝一角，米七萬六千三百三十二石九斗。}

\Method{以衰分求之。列母子求田率，副併為法，以共田為寔，寔如法而一，得一分之率。以徧乗未併者，得三等田。各以起租乗之，各得米。}

\Working{置弱半母四為中率，子三為上率。以太半子二减母三，餘一，以乗中率四，只得四為中泛。又以餘一乗上率三，只得三為上泛。次以太半母三乗中泛四，得一十二為下泛。副併三泛，得一十九為法。置田三千二十一頃五十一畆，以畝法二百四十通之，得七千二百五十一萬六千二百四十歩，内子一十五歩，得七千二百五十一萬六千二百五十五歩為寔。以法一十九除之，得三百八十一萬六千六百四十五歩為一分之數。以上泛三因之，得一千一百四十四萬九千九百三十五歩為上積。又以中泛四因之一分數，得一千五百二十六萬六千五百八十歩為中積。又以下泛一十二乗一分數，得四千五百七十九萬九千七百四十歩為下積。其三積各以畝法二百四十約之為畆。其上田得四百七十七頃八畝一十五歩，其中田得六百三十六頃一十畝三角，其下積得一千九百八頃三十二畝一角。各以起租，三積為三寔。其上積一千一百四十四萬九千九百三十五歩乗上積六斗，得六百八十六萬九千九百六十一石為寔，以畝法二百四十除之，得二萬八千六百二十四石八斗三升七合五勺為上田租。其中積一千五百二十六萬六千五百八十歩乗中田租四斗五升，得六百八十六萬九千九百六十一石為寔，以畝法二百四十除之，得二萬八千六百二十四石八斗三升七合五勺。其下積四千五百七十九萬九千七百四十歩乗下租四斗，得一千八百三十一萬九千八百九十六石為寔，以畝法二百四十除之，得七萬六千三百三十二石九斗為下田米。}

\EditorialNote{草内求各衰數，意謂以三分為上田衰數，加弱半三分之一得四分為中田衰數，三因中田衰數得一十二分為下田衰數，併之得一十九分為總衰數。其法本属顯明，但語中參入母子副泛等名目，其意反晦矣。}


%% =============================================================================
%% FASCICLE 5B — 賦役 (Part 2)
%% =============================================================================
\newpage
\section*{\color{titlecolor}卷五下\quad 賦役}
\addcontentsline{toc}{section}{卷五下\quad 賦役}

\subsection*{\color{sectioncolor}戶稅移割}
\addcontentsline{toc}{subsection}{戶稅移割}

\Question{某縣據甲稱：本户田地原納苗三十五石七斗，和買本色一十一匹二丈二尺九寸四分八釐七毫五絲，折帛二十七匹二寸一分三釐七毫五絲，紬絹折帛八匹三丈九寸七寸三分九釐，紬絹本色二十匹二丈八尺九寸三分九釐。已将田四百七畆出與乙，五百十六畆出與丙。訖乞移割本户所出田上税賦，歸併乙丙兩户送納。㑹到乙原有田三百七十五畆，丙原有田四百六十三畆，並係本鄉本等。毎畆苗三升五合，税一尺一寸五分，物力一貫二百；本等田紬一尺三寸四分，物力九百；物力及三十二貫數和買一匹，内三分本色七分折帛；夏税七分本色三分折帛；紬五分本色五分折帛，併絹紬。欲知甲田地及甲乙丙分割合納苗米、和買、夏税、折帛、本色、畸零各幾何？}

\EditorialNote{此題科税分田地二項。田有科米、税絹、物力三色，地只有税紬、物力二色，絹與紬折色不同，物力和買合田地總計之。又甲户有田者有地，乙丙二户有田無地。此等皆不可不明言者，題中殊未分晰。}

\Answer{甲原有田一千二十畆，地一十一畆二角四十八步。

甲今有田九十七畆，苗米三石三斗九升五合，夏税折帛三丈三尺四寸六分五釐，本色一匹畸零三丈八尺八分五釐，物力一百一十六貫四百文。地一十一畆二角四十八步，税紬折帛七尺八寸三分九釐，税紬畸零七尺八寸三分九釐，物力一十貫五百三十文，田地共物力一百二十六貫九百三十文。和買折帛二匹三丈一尺六分三釐七毫五絲，本色一疋畸零七尺五寸九分八釐七毫五絲。

乙今有田七百八十二畆，苗米二十七石三斗七升，夏税折帛六匹二丈九尺七寸九分，本色一十五匹畸零二丈九尺五寸一分，物力九百三十八貫四百文，和買折帛二十匹二丈一尺一寸，本色八匹畸零三丈一尺九寸。

丙今有田九百七十九畆，苗米三十四石二斗六升五合，夏税折帛八匹一丈七尺七寸五分五釐，本色一十九匹畸零二丈八尺九分五釐，物力一千一百七十四貫八百文，和買折帛二十五匹二丈七尺九寸五分，本色一十一匹畸零五寸五分。}

\Method{以粟米及衰分求之。置甲原納米為實，以毎畆苗為法除之，得甲原有田。以乗毎畆税，得田絹。次以甲紬絹本色折帛併之，内減田絹，餘為地紬。以毎紬約之，得甲原有地。次以甲出田併乙丙原有田，各得三户今有田地。各以等則乗之，各得物力苗税。次以毎畆物力率約各户共物力，得和買。副之，三分因之退位為和本，七分因之退位為和折；夏税反其分而因之退之；紬以半之，併入税，各得。}

\subsection*{\color{sectioncolor}均科綿税}
\addcontentsline{toc}{subsection}{均科綿税}

\Question{縣科綿有五等户，共一萬一千三十三户，共科綿八萬八千三百三十七兩六錢。上等一十二户，副等八十七户，中等四百六十四户，次等二千三十五户，下等八千四百三十五户。欲令上三等折半差，下二等此中等六四折差。科率求之，問各户納及各等幾何？}

\EditorialNote{題言下二等比中等六四折差，是次等為中等六分之四，下等又為次等六分之四也。乃術草中皆以十分之四收之，是四分折差而非四六折差矣。集中題與術不相應者多類此，豈成書之後未能重加校正歟？}

\Answer{上等一户一百二十四兩，一十二户計一千四百八十八兩；
副等一户六十二兩，八十七户計五千三百九十四兩；
中等一户三十一兩，四百六十四户計一萬四千三百八十四兩；
次等一户一十二兩四錢，二千三十五户計二萬五千二百三十四兩；
下等一户四兩九錢六分，八千四百三十五户計四萬一千八百三十七兩六錢。}

\Method{列五等户數，先以四折下等數，加次等户，又以四折之，加中等户數，却以半折之，加二等户，又以半折之，加上等户數，不折便為法。除科綿，得上等一户之綿。復半之為副等，又半之為中等，又四折之為次等，又四折之為下等。各一户綿，却各以户數乗之，各得五等共出綿。}

\Working{先置下等户八千四百三十五，以四折之，得三千三百七十四。乃以得數折入次户二千三十五内，得五千四百九户訖。又以四折次數五千四百九户，得二千一百六十三户六分。乃以得次數併入中户四百六十四内，共得二千六百二十七户六分。次以五分折中數二千六百二十七户六分，得數。次以得數一千三百一十三户八分併副户八十七，得一千四百户八分。又以五折副數一千四百户八分，得七百户四分，既得數。乃以副數七百户四分併上户一十二户，得七百一十二户四分，不折便為法。

累折至等，共得七百一十二户四分以為總法。乃置科綿八萬八千三百三十七兩六錢為實，法七百一十二户四分而一，得一百二十四兩為上等一户所出綿。以五折之，得六十二兩為副等一户所出綿。又五折之，得三十一兩為中等一户所出綿。乃以四折之，得一十二兩四錢為次等一户所出綿。又以四折之，得四兩九錢六分為下等一户所出綿。乃以各等户數各乗一户所出，即各得每等共數之綿。}

\subsection*{\color{sectioncolor}均定勸分}
\addcontentsline{toc}{subsection}{均定勸分}

\Question{欲勸糶賑濟，據甲民户物力畆步排定，共計一百六十二户，作九等：上等三户，第二等五户，第三等七户，第四等八户，第五等十三户，第六等二十一户，第七等二十六户，第八等三十四户，第九等四十五户。今先觀諭，第一等上户願糶五千石，第九等户願糶二百石。欲知各等抛差石數併數認米數各幾何？}

\Answer{認該米二十三萬七千六百石。

上等一户米五千石，三户計一萬五千石；
二等一户米四千四百石，五户計二萬二千石；
三等一户米三千八百石，七户計二萬六千六百石；
四等一户米三千二百石，八户計二萬五千六百石；
五等一户米二千六百石，一十三户計三萬三千八百石；
六等一户米二千石，二十一户計四萬二千石；
七等一户米一千四百石，二十六户計三萬六千四百石；
八等一户米八百石，三十四户計二萬七千二百石；
九等一户米二百石，四十五户計九千石。}

\Method{以衰分求之。置上下户米，減餘為實。列等數減一，餘為法。除之，得抛差石數。以差累減上等米，各得諸等米。以各等户乗之，併之為總數。}

\Working{置上等户米五千石，減下等户二百石，餘四千八百石為實。以九等減一，餘八為法。除實，得六百石為每等抛差。用減上等，餘四千四百石為二等米。又減六百，得三千八百石為三等米。又減六百，得三千二百石為四等米。又減六百，得二千六百石為五等米。又減六百，得二千石為六等米。又減六百，得一千四百石為七等米。又減六百，得八百石為八等米。又減六百，得二百石為九等米。又各乗户數，併之，得總認米石數二十三萬七千六百石。}

\subsection*{\color{sectioncolor}均定合解}
\addcontentsline{toc}{subsection}{均定合解}

\Question{州郡合解諸司窠名錢：户部九十六萬五千四百二十一貫文，總所六十四萬三千六百一十四貫文，運司一萬六千九十貫三百五十文。今諸窠名先催到九千二百五十三貫六百二十文，欲照原額分數均定樁收解，合各幾何？}

\Answer{户部五千四百九十七貫二百文；總所三千六百六十四貫八百文；運司九十一貫六百二十文。}

\Method{以衰分求之。置諸原率，可約約之。副併為法。以催到錢乗未併者，各為實。實如法而一。}

\Working{列户部九十六萬五千四百二十一貫，總所六十四萬三千六百一十四貫，運司一萬六千九十貫三百五十文，各為原率。今原率可約，求等得一萬六千九十貫三百五十為等數，俱約之：户部得六十，總所得四十，運司得一，各為率。副併得一百一為法。以催到錢九千二百五十三貫六百二十文乗未併數：六十、四十及一，畢得五十五萬五千二百一十七貫二百文為户部之實，三十七萬一百四十四貫八百文為總所之實，九千二百五十三貫六百二十文為運司之實。三實皆如法一百一而一：其户部得五千四百九十七貫二百文，總所得三千六百六十四貫八百文，運司得九十一貫六百二十，各為候解錢。}

\newpage
\subsection*{\color{sectioncolor}寬減屯租}
\addcontentsline{toc}{subsection}{寬減屯租}

\Question{屯租欲議寬減，仍聴以夏麥折納分數。官牛種者與減二分，私牛種者與減四分。每嵗租榖以三分之一許夏折二麥，内四分大六分小折色。每大麥三石折小麥二石，小麥二石折榖三石五斗。屯租舊額官種一石納租五石，私種一石納租三石。今某州屯田去年計官私種共九千七百八十二石，共合收租榖三萬九千五百八十六石。欲知官私種各數、原額、今減、合催、成年夏麥秋榖幾何？}

\EditorialNote{此題有官私共種數、租數，有種一石租數，求各種數租數，古謂之差分，今謂之和較比例。折色亦差分也，今謂之和數比例。大小麥與榖相易，古謂之異乗同乗，今謂之和率比列。}

\Answer{官種五千一百二十石，私出種四千六百六十二石。

原租額共三萬九千五百八十六石：官種二萬五千六百石，私種一萬三千九百八十六石。

今減一萬七百一十四石四斗：官種五千一百二十石，私種五千五百九十四石四斗。

合催二萬八千八百七十一石六斗。

官種租二萬四百八十石：一分折麥計榖六千八百二十六石六斗六升零三分升之二；四分折大麥二千三百四十石五斗七升七分升之一（係折榖二千七百三十石六斗六升三分升之二）；六分折小麥二千三百四十石五斗七升七分升之一（係折榖四千九十六石二分）；正色榖一萬三千六百五十三石三斗三升三分升之一。

私種租八千三百九十一石六斗：一分折麥計榖二千七百九十七石二斗；四分折大麥九百五十九石四升（係折榖一千一百一十八石八斗八升）；六分折小麥九百五十九石四升（係折榖一千六百七十八石三斗二升）；二分正色榖五千五百九十四石四斗。

成年共收：夏折榖九千六百二十三石八斗六升三分升之二；大麥三千二百九十九石六斗一升七分升之一；小麥三千二百九十九石六斗一升七分升之一；秋正榖一萬九千二百四十七石七斗三升三分升之一。}

\subsection*{\color{sectioncolor}戶稅均寬}
\addcontentsline{toc}{subsection}{戶稅均寬}

\Question{州郡寬恤，近将某縣下三等税户秋科餘欠錢米，已與蠲放：共錢一千三百五十五貫七百六文，米五千二百七十二石一斗九升。某本縣下三等物力計三萬七千六百五十八貫五百文。今來官員陳述：本縣多有樂輸無欠之户，今䝉蠲放税尾，似反寬潤頑輸之户，於理未均。遂議将樂輸三等户於明年兩税與照昨來體例減免。契勘得三等無欠户物力二十二萬八千一十五貫三百二十一文。欲知每百文合減免錢米及共減各幾何？}

\Answer{毎物力一百文，放錢三文六分，放米一升四合。明年兩税放免：錢七千九百四十九貫三百五十一文五分五釐六毫；米三萬九百一十四石一斗四升四合九勺四抄。}

\Method{以粟米衰分求之。置原物力為法，除原放錢米，得毎百文物力所放錢米率。以各率乗今來物力，各得錢米，為明年兩税合放數。}

\Working{以原物力三萬七千六百五十八貫五百文為法。先除原放錢一千三百五十五貫七百六文，定法百文上得一文為商，除得三文六分為物力百文所放錢率。次除原放米五千二百七十二石一斗九升，得一升四合為物力百文所放米率。以今三等户物力二十二萬八千一十五貫三百二十一文徧乗所放錢米二率，得錢七千九百四十九貫三百五十一文五分五釐六毫，得米三萬九百一十四石一斗四升四合九勺四抄，為明年兩税合放數。}

\subsection*{\color{sectioncolor}米榖粒分}
\addcontentsline{toc}{subsection}{米榖粒分}

\Question{開倉受約，有甲户米一千五百三十四石。到廊驗得米内夾榖，乃於様内取米一捻，數計二百五十四粒，内有榖二十八顆。凡粒米率每勺三百。今欲知米内雜榖多少，以折米數科貴及粒各幾何？}

\Answer{米一千三百六十四石八斗九升七合六勺一百二十七分勺之四十八；
榖一百六十九石一斗二合三勺一百二十七分勺之七十九。
合折米八十四石五斗五升一合一勺一百二十七分勺之一百三。
原米折米共計四十三億四千八百三十四萬六千四百五十六粒。}

\Method{以粟米求之，衰分入之。置様米粒數為法，以常榖顆數減之，餘與榖為列衰。可約約之。以共米乗列衰為各實，實如法而一，各得米數榖數。置榖數以粟率折之，為榖所折米。次以勺率遍乗米數折米，得粒數。}

\Working{置一捻様粒數二百五十四為法，以帶榖二十八顆為榖衰，以減法，餘二百二十六為米衰。此二衰與法皆可約，求等得二，俱以約之：法得一百二十七，米衰得一百一十三，榖衰得一十四。以米一千五百三十四石遍乗二衰，得一十七萬三千三百四十二石為米實，得二萬一千四百七十六石為榖實。皆如法一百二十七而一：米得一千三百六十四石八斗九升七合六勺一百二十七分勺之四十八，榖得一百六十九石一斗二合三勺一百二十七分勺之七十九。以粟率五十折之，得八十四石五斗五升一合一勺一百二十七分勺之一百三為榖折納米數。併二米，得一千四百四十九石四斗四升八合八勺一百二十七分勺之二十四。先通分納子一十八萬四千八十石，以勺率三百粒乗之，得五千五百二十二億四千萬粒為實，以母一百二十七除之，得四十三億四千八百三十四萬六千四百五十六粒，不盡八十八棄之。}


%% =============================================================================
%% FASCICLE 6A — 錢穀 (Part 1)
%% =============================================================================
\newpage
\section*{\color{titlecolor}卷六上\quad 錢穀}
\addcontentsline{toc}{section}{卷六上\quad 錢穀}

\subsection*{\color{sectioncolor}課糴}
\addcontentsline{toc}{subsection}{課糴}

\Question{差人五路和糴。據浙西平江府石價三十五貫文，一百三十五合，至鎮江水脚錢每石九百文；安吉州石價二十九貫五百文，一百一十合，至鎮江水脚錢每石一貫二百文；江西隆興府石價二十八貫一百文，一百一十五合，至建康水脚錢每石一貫七百文；吉州石價二十五貫八百五十文，一百二十合，至建康水脚錢每石二貫九百文；湖南潭州石價二十七貫三百文，一百一十八合，至鄂州水脚錢每石一貫七百文。其錢並十七界官㑹，其米並用文思院斛交量紐數。欲皆以官解計石錢相比貴賤幾何？}

\EditorialNote{按草中係二貫一百文，此訛。文思院解每斗八十三合。}

\Answer{文思院斛石錢：安吉州二十三貫一百六十四文一十一分文之六；平江府二十二貫七十一文二十七分文之二十三；隆興府二十一貫五百七文二十三分文之一十九；潭州二十貫六百七十九文五十九分文之四十九；吉州一十九貫八百八十五文十二分文之五。}

\EditorialNote{按三十九訛四十九。}

\Method{以粟米互換求之。置石價併水脚，乗石數，又乗官斗合數為實。各如本州合數而一，各得官解石錢。以課貴賤。}

\Working{置安吉州石價二十九貫五百文，平江石價三十五貫文，隆興石價二十八貫一百文，吉州石價二十五貫八百五十文，潭州石價二十七貫三百文，列右行。次置水脚：安吉一貫二百文，平江九百文，隆興一貫七百文，吉州二貫九百文，潭州二貫一百文，列左行，各對本州石價。以兩行數併之，得數：安吉三十貫七百文，平江三十五貫九百，隆興二十九貫八百，潭州二十九貫四百，吉州二十八貫七百五十，仍於右行。次以文思院官斗八十三合遍乗之：安吉州得二千五百四十八貫一百文，平江府得二千九百七十九貫七百文，江西隆興得二千四百七十三貫四百文，湖南潭州得二千四百四十貫二百文，江南吉州得二千三百八十六貫二百五十文，各為實於右。得次列安吉斗一百一十合，平江斗一百三十五合，隆興斗一百一十五合，潭州斗一百一十八合，吉州斗一百二十合於左行為法。以对除右行之實：安吉得二十三貫一百六十四文一十一分之六，平江得二十三貫七十一文二十七分文之二十三，隆興得二十一貫五百七文二十三分文之一十九，潭州得二十貫六百七十九文五十九文分之四十九，吉州得一十九貫八百八十五文一十二分文之五。相課石價，其安吉州最貴，平江次之，隆興又次之，潭州又次之，吉州最賤。}

\newpage
\subsection*{\color{sectioncolor}折解輕齎}
\addcontentsline{toc}{subsection}{折解輕齎}

\Question{有甲乙丙丁四郡，各合起上供銀絹。

\textbf{甲郡}：銀三千二百兩，每兩二貫二百文足；絹六萬四千匹，每匹二貫文足。去京一千里，每擔一里傭錢六文足。其時舊㑹每貫五十四文足。

\textbf{乙郡}：銀二千七百兩，每兩二貫三百文足；絹四萬九千二百匹，每匹二貫四百二十文足。去京九百八十里，每擔一里傭錢四文二分足。舊㑹價五十九文足。

\textbf{丙郡}：銀四千兩，每兩新㑹九貫三百文；絹七萬三千六百匹，每匹新㑹一十貫三百文。去京二千里，每擔一里傭錢八十文舊㑹。

\textbf{丁郡}：銀二千六百兩，每兩五十一貫文舊㑹；絹三萬二千三十五匹，每匹五十八貫文舊㑹。去京一千五百里，每擔一里傭錢一百文舊㑹。

諸郡銀每五百兩、絹每六十匹、新㑹每五千貫為擔。欲並折新㑹，均作三限起解。求各郡每限及本色原理、折解實用、寛餘、傭錢各新㑹幾何？}

\EditorialNote{此題貫數分三項：其一足數，每千文為一貫；其一舊㑹數，如甲以五十四文為一貫，乙以五十九文為一貫是也；其一新㑹數，為舊㑹數之五倍，如甲以二百七十文為一貫，乙以二百九十五文為一貫是也。四郡或言足數，或言舊㑹數、新㑹數。並折新㑹數。傭錢原以銀五百兩或絹六十匹各為一擔，今皆折新㑹數，以五千貫為一擔，故有寛餘錢數。題語多未詳，而併為一擔句尤混，略為分析而術草之意大概可見矣。}

\Answer{
\textbf{甲郡}：合解五十萬一百四十八貫一百四十八文。初限一十六萬六千七百一十六貫四十九文，次限一十六萬六千七百一十六貫四十九文，未限一十六萬六千七百一十六貫五十文。傭錢原理二萬三千八百四十五貫九百二十五文二十七分文之二十五，實用二千二百二十二貫八百七十七文二十七分文之二十一，寛餘二萬一千六百二十三貫四十八文二十七分文之四。

\textbf{乙郡}：合解四十二萬四千六百五十七貫六百二十七文。初限一十四萬一千五百五十二貫五百四十二文，次限一十四萬一千五百五十二貫五百四十二文，末限一十四萬一千五百五十二貫五百四十三文。傭錢原理一萬一千五百一十六貫四百二十八文五十九分文之二十八，實用一千一百八十五貫一十文二百九十五分文之五，寛餘一萬三百三十一貫四百一十八文五十九分文之二十七。

\textbf{丙郡}：合解七十九萬五千二百八十貫文。初限二十六萬五千九十三貫三百三十三文，次限二十六萬五千九十三貫三百三十三文，末限二十六萬五千九十三貫三百三十四文。傭錢原理三萬九千五百九貫三百三十三文三分文之一，實用五千八十九貫七百九十二文，寛餘三萬四千四百一十九貫五百四十一文三分文之一。

\textbf{丁郡}：合解三十九萬八千一百二十六貫文。初限一十三萬二千七百八貫六百六十六文，次限一十三萬二千七百八貫六百六十六文，末限一十三萬二千七百八貫六百六十八文。傭錢原理一萬四千一百七十三貫五百文，實用二千三百八十八貫七百五十六文，寛餘一萬一千七百八十四貫七百四十四文。}

\Method{以均輸求之。置各郡銀絹，乗各價，併之，歸足原，展足為舊㑹，次以五約舊㑹為新㑹，各得合解錢。以限數除之，得每限錢，不盡併歸末限。次置里數，乗每里傭價為率。以率乗原銀及原絹，各為傭實。以每擔銀絹率各為法，實如法而一，不滿者亦為擔，併之為原理傭錢。次以率乗合觧錢為實，乃以錢物毎擔率為法，實如法而一，各得實用傭錢。以減原理傭錢，餘為寛餘傭錢。}


%% =============================================================================
%% FASCICLE 6B — 錢穀 (Part 2)
%% =============================================================================
\newpage
\section*{\color{titlecolor}卷六下\quad 錢穀}
\addcontentsline{toc}{section}{卷六下\quad 錢穀}

\subsection*{\color{sectioncolor}僦直推原}
\addcontentsline{toc}{subsection}{僦直推原}

\Question{房廊數内，一户日納一百五十六文八分足為準。指揮未曽經減者，減三分；已曽經減三分者，減二分；已曽經減二分者，更減二分。今本户累經減者，欲知原額房錢幾何？}

\Answer{原額三百五十文。}

\Method{以衰分求之。列一十分兩行，各三位；列減分，對減右行。以餘者相乘為法。以左行原列相乘，得納錢為實。實如法而一，得原額錢。}

\Working{列一十三位於左行，又列一十分三位於右行。其右上減去初減三分，右中減去次減二分，右下減去更減二分。右行餘七八八，以相乗得四百四十八為法。乃以左行三位一十分相乘，得一千為乗率。以乘見日納錢一百五十六文八分，得一百五十六貫八百文為實。實如法而一，得三百五十文，為本户原額戸錢。}

\subsection*{\color{sectioncolor}推求本息}
\addcontentsline{toc}{subsection}{推求本息}

\Question{三庫息例：萬貫以上一釐，千貫以上二釐五毫，百貫以上三釐。甲庫本四十九萬三千八百貫，乙庫本三十七萬三百貫，丙庫本二十四萬六千八百貫。今三庫共約到息錢二萬五千六百四十四貫二百文。其典率：甲反錐差，乙方錐差，丙蒺藜差。欲知原典三例本息各幾何？}

\EditorialNote{此即衰分題也。其差有反錐、方錐、蒺藜之名，蓋以一二三遞減如立錐為反錐，以一四九平方遞加為方錐，以一三六三數遞加為蒺藜。是必古有其名也。至以各差求各本，則因各本原依各差入之也。}

\Answer{\textbf{甲庫}共納息九千五十三貫文：一釐息二千四百六十九貫文，二釐半息四千一百一十五貫文，三釐息二千四百六十九貫文。

\textbf{乙庫}共納息一萬五十一貫文：一釐息二百六十四貫五百文，二釐半息二千六百四十五貫文，三釐息七千一百四十一貫五百文。

\textbf{丙庫}共納息六千五百四十貫二百文：一釐息二百四十六貫八百文，二釐半息一千八百五十一貫文，三釐息四千四百四十二貫四百文。}

\Method{置諸庫諸色之差，照釐率為三行。縱併之為約率，横命之為乗率。以約率各約自庫之本，各得。以遍乗未併乗率，然後各以釐率横乗之，次以縱併之，為各庫共息。}

\Working{置甲庫反錐差，自下置三二一於右行。次置乙庫方錐差，自上置一四九於中行。次置丙庫蒺藜差，自上置一三六於左行。各為三庫上中下三等乘率。乃縱併甲差三二一，得六為甲約率。縱併乙差一四九，得一十四為乙約率。縱併丙差一三六，得一十為丙約率。直命九位數各為上中下乗率。乃先以約率各約自庫之本：乃以甲約率六約甲本四十九萬三千八百貫，得八萬二千三百貫為甲得。次以乙約率一十四約乙本三十七萬三百貫，得二萬六千四百五十貫為乙得。次以丙約率一十約丙本二十四萬六千八百貫，得二萬四千六百八十貫為丙得。以各得乗未併乗率：其甲所得八萬二千三百貫，乘反錐乗率三二一，得二十四萬六千九百貫為上率，得一十六萬四千六百貫為中率，得八萬二千三百貫為下率。其乙所得二萬六千四百五十貫，以乗方錐差一四九，得二萬六千四百五十貫為上率，得一十萬五千八百貫為中率，得二十三萬八千五十貫為下率。其丙所得二萬四千六百八十貫，以乗蒺藜差一三六，得二萬四千六百八十貫為上率，得七萬四千四十貫為中率，得一十四萬八千八十貫為下率。然後各以息釐數乘各庫三乗：其甲以一釐乗上率二十四萬六千九百貫，得二千四百六十九貫為上息；以二釐五毫乘中率一十六萬四千六百貫，得四千一百一十五貫為中息；以三釐乘下率八萬二千三百貫，得二千四百六十九貫為下息。併上中下三息，得九千五十三貫文為甲庫共息。其乙庫以一釐乗上率二萬六千四百五十貫，得二百六十四貫五百文為上息；以二釐五毫乗中率一十萬五千八百貫，得二千六百四十五貫為中息；以三釐乗下率二十三萬八千五十貫，得七千一百四十一貫五百為下息。併上中下三息，得一萬五十一貫文為乙庫共息。其丙庫以一釐乗上率二萬四千六百八十貫，得二百四十六貫八百為上息；以二釐五毫乘中率七萬四千四十貫，得一千八百五十一貫為中息；以三釐乗下率一十四萬八千八十貫，得四千四百四十二貫四百文為下息。併上中下三息，得六千五百四十貫二百文為丙庫共息。併三庫共息，得二萬五千六百四十四貫二百文為總息。}

\subsection*{\color{sectioncolor}易牒知原}
\addcontentsline{toc}{subsection}{易牒知原}

\EditorialNote{按舊本此問無題，今増入。}

\Question{出度牒差人營運：毎三道易鹽一十三袋，鹽二袋易布八十四疋，布一十五疋易絹三疋半，絹六疋易銀七兩二錢。今趂到銀九千一百七十二兩八錢。欲知原闗度牒道數幾何？}

\Answer{度牒一百八十道。}

\Method{以粟米互乘易法求之。列各數以本色相對，如鴈翅。以多一事者相乗為實，以少一事者相乘為法，除之。}

\Working{先以度牒三道乗鹽二袋，得六。以乘布一十五，得九十。又乗絹六疋，得五百四十。乃乗銀九萬一千七百二十八錢，得四千九百五十三萬三千一百二十錢為實。次以鹽一十三袋乗布八十四，得一千九十二。以乗絹三疋五分，得三千八百二十二。乃乘銀七兩二錢，得二十七萬五千一百八十四錢為法。除實，得一百八十道為原闗度牒。}

\subsection*{\color{sectioncolor}粟米交易}
\addcontentsline{toc}{subsection}{粟米交易}

\EditorialNote{按舊本此問無題，今増。}

\Question{菽三升易小麥二升，小麥一升五合易油麻八合，油麻一升二合易粳米一升八合。今將菽一十四石四斗欲易油麻，又將小麥二十一石六斗欲易粳米。問各幾何？}

\Answer{油麻五石一斗二升；粳米一十七石二斗八升。}

\Method{以粟米換易求之。置原易率，本色對列，如鴈翅。以多一事者相乘為實，以少一事者相乗為法，除之。各得或問數，不干其率者不置。}

\Working{置四色六數，列六位率，如鴈翅，皆化為合。先將菽一十四石四斗化作一萬四千四百合，乃對前二句率數四位，如鴈翅，至欲易油麻止，共五事為上圖。次將小麥二十一石六斗化為二萬一千六百合，乃對後兩句率四位，鴈翅，至欲粳米止，共五事為下圖。

下圖：其上圖以菽一萬四千四百合乘麥二十，得二十八萬八千，又乗油麻八合，得二百三十萬四千合為油麻實。次以菽三十合乘麥一十五合，得四百五十合為法。除之，得五千一百二十合，展為五石一斗二升為油麻。

其下圖以小麥二萬一千六百合乗油麻八合，得一十七萬二千八百合，又乗粳米一十八合，得三百一十一萬四百合為粳米實。以小麥一升五合乗油麻一十二合，得一百八十合為法。除之，得一萬七千二百八十，展作一十七石二斗八升為粳米。}

\subsection*{\color{sectioncolor}計米易麴}
\addcontentsline{toc}{subsection}{計米易麴}

\EditorialNote{按舊本此問無題，今増。}

\Question{庫率：粳榖七石出米三石，糯米一斗易小麥一斗七升，小麥五升踏麴二斤四兩，麴一百一十斤醞糯米一石三斗。今有糯穀一千七百五十九石三斗八升，欲出穀做米，易麥踏麴，還自醞，餘穀之米須令適足。各合幾何？}

\Answer{共穀一千七百五十九石三斗八升。出穀九百二十四石，得米三百九十六石。易麥六百七十三石二斗。踏麴三萬二百九十四斤。餘榖八百三十五石三斗八升。醞米三百五十八石二升。}

\Method{以粟米換易求之。置諸率，隨本色對列，如鴈翅。有分者通之，異類者變之，以頭位者進乗之，以下位退乗之，得合數。有對者相乗之，無對者直命之，為諸率。併上下無對者為法率。以今有物徧乘諸率（不乗法率），各為實。諸實並如法而一，各得其已變者，復互易乗除之，即得所求。}

\Working{置糯穀七、出米三於右行上副兩位。次置糯米一斗、麥一斗七升於副行副中兩位。次置小麥五升、踏麥二斤四兩於次行中次兩位。次置麴一百一十觔、醞米一石三斗於左行次下兩位，隨本色對列，如鴈翅。訖乃驗次行二斤四兩是四分斤之一，以母四通次行兩位，以子一内次行次位，具中位得二，次位得九。又驗左行下位是糯米，是異類，於糯米合變為糯穀，乃以問中首句率穀七米三變之，以七因米一石三斗得九石一斗於左下為穀，却以米三因麴一百一十斤為三百三十斤麴於左行，得變圖數。以左行三百三十乗次行二，得六百六十。次以六百六十乘副行一十，得六千六百。次以六千六百乗右上七，得四萬六千二百，各於原位。却以右行副位三因副行一斗七升，得五斗一升。又以五斗一升乗次行九，得四百五十九。又以四百五十九乗左下九十一，得四萬一千七百六十九，列為合圖數。乃驗合圖四行，其副中次三位有對，以對相乗合之。其右上左下無對者，直命之，皆為率。列右行：上得四萬六千二百為出糯穀率，副位得一萬九千八百為得糯米率，中得三萬三千六百六十為易得麥率，次得一萬五千一十四七為踏到麴率，下得四萬一千七百六十九為餘下糯榖率。併上下率，共得八萬七千九百六十九為法率。今六率共求等得一，約之，只得原率為率圖。始用今有穀一千七百五十九石三斗八升，皆化為升，徧乘五率（不乗法率），得八十一億二千八百三十三萬五千六百升為出穀實，得三十四億八千三百五十七萬二千四百升為糯米實，得五十九億二千二百七萬三千八十升為易麥實，得二十六億六千四百九十三萬二千八百八十六為踏麴實，得七十三億四千八百七十五萬四千三百二十二升為餘穀實。其五實皆如法八萬七千九百六十九而一：得九百二十四石為出穀，得三百九十六石為做到糯米，得六百七十三石二斗為易到小麥，得三萬二百九十四斤為踏到麴，得八百三十五石三斗八升為餘下穀。今將餘下穀變為米，乃以乗率三因餘穀八百三十五石三斗八升，得二千五百六石一斗四升為實，以糯穀率七為法除之，得三百五十八石二升為醞米。}

\subsection*{\color{sectioncolor}囘運費}
\addcontentsline{toc}{subsection}{囘運費}

\Question{有江西水運米一十二萬三千四百石，原係鎮江交卸，計水程二千一百三十里，每石水脚錢一貫二百文。今截上件米就池州安頓，池州至鎮江八百八十里。欲收囘不該水脚錢幾何？}

\Answer{收囘錢六萬一千一百七十八貫五百九十一文。}

\Method{以粟米互易求之。置池州至鎮江里數，乗水脚錢，得數又乗運米為實。以原至鎮江水程為法，除實，得收囘錢。}

\Working{置池州至鎮江八百八十里，乗毎石水脚錢一貫二百，得一千五十六貫文。又乗運米一十二萬三千四百石，得一億三千三十一萬四百貫文為實。以原至鎮江水程二千一百三十里為法，除實，得六萬一千一百七十八貫五百九十一文為收囘錢。}

\subsection*{\color{sectioncolor}三色均價}
\addcontentsline{toc}{subsection}{三色均價}

\Question{庫有三色金共五千兩：内八分金一千二百五十兩，兩價四百貫文；七分五釐金一千六百兩，兩價三百七十五貫文；八分五釐金二千一百五十兩，兩價四百二十五貫文。並欲煉為足色，每兩工食藥炭錢三貫文，耗金九百七十二兩五錢。欲知色分及兩價各幾何？}

\Answer{色十分；兩價五百三貫七百二十四文，五百三十七分文之二百一十二。}

\Method{以方田及粟米求之。置共數，以耗減之，餘為法。以三色分數各乗兩數，併之為色分實。以三色價數各乗兩數為寄，以工藥價乗共金，併寄，共為價實。二實皆如法而一，即各得。}

\Working{置共金五千兩，減耗九百七十二兩五錢，外餘四千二十七兩五錢為法。次置一千二百五十兩，乘八分，得一萬分於上。置一千六百兩，以七分五釐乗之，得一萬二千分，加上。置二千一百五十兩，乗八分五釐，得一萬八千二百七十五分，又加上，共得四萬二百七十五分為分實。次置一千二百五十兩乗四百貫，得五十萬貫為寄。次置一千六百兩乘價三百七十五貫，得六十萬貫，加寄。次置二千一百五十兩乗價四百二十五貫，得九十一萬三千七百五十貫，又加寄。次置共金五千兩，乗工藥錢三貫，得一萬五千貫，又加寄，共得二百二萬八千七百五十貫為貫實。二實並如法四千二十七兩五錢而一：得其色一十分；其價每兩得五百三貫七百二十四文，不盡一貫五百九十。與法求等，得七十五，俱約之，為五百三十七文之二百一十二。}


%% =============================================================================
%% END MATTER
%% =============================================================================
\newpage
\section*{\color{titlecolor}Editorial Notes}
\addcontentsline{toc}{section}{Editorial Notes}

This digital edition of Qin Jiushao's \textit{Shuxue Jiuzhang} (數學九章) Fascicles 5--6 presents the complete text of the ``Taxation and Corv\'{e}e Labor'' (賦役) and ``Money and Grain'' (錢穀) chapters, transcribed from the \textit{Siku Quanshu} (四庫全書) edition as preserved on Chinese Wikisource.

\subsection*{Structure of the Text}

Each problem in the \textit{Shuxue Jiuzhang} follows a traditional four-part structure:

\begin{enumerate}
    \item \textbf{問 (W\`{e}n)} --- The problem statement, presenting the practical scenario and given quantities.
    \item \textbf{答 (D\'{a})} --- The answer, listing all required numerical results.
    \item \textbf{術 (Sh\`{u})} --- The method, describing the mathematical approach in general terms.
    \item \textbf{草 (C\v{a}o)} --- The working, showing the detailed step-by-step computation.
\end{enumerate}

\subsection*{Historical Context}

Qin Jiushao (c.\ 1202--1261), style name Daogu (道古), was one of the foremost mathematicians of the Song Dynasty. His \textit{Shuxue Jiuzhang}, completed in 1247, is divided into nine categories:

\begin{enumerate}
    \item \textbf{Dayan} (大衍) --- Indeterminate analysis
    \item \textbf{Tianshi} (天時) --- Calendrical computations
    \item \textbf{Tianyu} (田域) --- Land measurement
    \item \textbf{Cewang} (測望) --- Surveying
    \item \textbf{Fuyi} (賦役) --- Taxation and corv\'{e}e labor
    \item \textbf{Qianu} (錢穀) --- Money and grain
    \item \textbf{Yingjian} (營建) --- Construction
    \item \textbf{Junlv} (軍旅) --- Military logistics
    \item \textbf{Shiyi} (市易) --- Trade and commerce
\end{enumerate}

\subsection*{Editorial Conventions}

The notes marked with ［按］(\`{A}n) are editorial annotations found in the \textit{Siku Quanshu} edition, often pointing out errors in the original text or discrepancies between the problem statement and the computational working. These annotations were added by the Qing Dynasty editors who collated the text for the imperial library.

\vspace{2cm}
\begin{center}
\textcolor{rulecolor}{\rule{0.5\textwidth}{0.4pt}}

\vspace{1em}
{\small\color{sectioncolor}
\textit{Typeset in XeLaTeX with Noto Sans CJK SC}\\
Traditional Chinese mathematical text structure preserved\\
Public Domain --- \textit{Siku Quanshu} Edition}
\end{center}

\end{document}
